\section{Setup and notation}
\label{sec:setup}

We keep the reference paper's model, recursions and symbol set, and add only
what Sections~\ref{sec:repro}--\ref{sec:nofactor} need. The plant is
$x_{k+1}=Ax_k+Bu_k+w_k$ with virtual-sensor reading $z_k=Fx_k+v_k$
($w_k,v_k$ i.i.d.\ Gaussian, $W\succ0$), and the reference cost is the average
stage cost of the LQG pair. Under the certainty-equivalence structure the control
law is $u_k=-K\hat x_k$ and the cost separates as
\begin{equation}
 D(S)\;=\;j_c+\operatorname{tr}(\Theta P_m),\qquad
 j_c=\operatorname{tr}(WP_c),
\label{eq:cost}
\end{equation}
with $P_c$ the control ARE solution, $\Theta=K^{\top}(R+B^{\top}P_cB)K$ the
closed-loop weighting, and $P_m$ the posterior error covariance. The
posterior is propagated in \emph{information form},
\begin{equation}
 P_m^{\,+}=\big(\tilde P^{\,-1}+S\big)^{-1},\qquad
 \tilde P^{\,+}=AP_mA^{\top}+W,
\label{eq:recursion}
\end{equation}
where $S\succeq0$ is the information matrix the virtual sensor injects per step,
and the rate is
\begin{equation}
 I(S)=\tfrac12\log_2\frac{\det(AP_mA^{\top}+W)}{\det P_m}
\label{eq:rate}
\end{equation}
evaluated at a stationary point of \eqref{eq:recursion}. A stationary point
exists exactly when the pair seen by the recursion is detectable, which is why
\eqref{eq:recursion} doubles as the feasibility certificate in
Proposition~\ref{pr:region}.

\emph{Admissibility.} A sensor is task-restricted to $F$ when
$\mathcal K_F=\{S\succeq0:\operatorname{range}(S)\subseteq
\operatorname{range}(F^{\top})\}$, and the paper's Corollary~1 parameterises the
cone as $S=U_FMU_F^{\top}$ with $M\succeq0$ and $U_F$ an orthonormal basis of
$\operatorname{range}(F^{\top})$. Two consequences are used here: the cone is
\emph{closed}, so rank-deficient $S$ --- a channel switched off entirely --- are
admissible, and the constraint is an equality-type reachability condition on a
subspace, not an inequality on a spectrum.

\emph{Notation.} $V:=\operatorname{range}(S)$, $P_V$ the orthogonal projector
onto $V$, $r:=\dim V$. We write $D_{\min}(V)$, not a function of $F$, because
every quantity in Sections~\ref{sec:infeasible}--\ref{sec:nofactor} depends on
$F$ only through the subspaces $F$ admits; this also avoids the symbol clash
noted in. Finally, writing $S=vP_V$ for the isotropic
family, the two evaluation calibres of
Remark~\ref{rmk:calibre} are
\[
 D^{\rm cone}_{V}(I)=\inf_{\substack{\operatorname{range}S\subseteq V\\
 I(S)\le I}}D(S),\qquad
 D^{\rm ray}_{V}(I)=\inf_{\substack{S=vP_V\\ I(S)\le I}}D(S),
\]
so $D^{\rm ray}\ge D^{\rm cone}$ on the same axis. Both are rate-\emph{budget}
calibres: $D$ is non-increasing and $I$ non-decreasing in the L\"owner order on
$S$, so the infimum sits on the boundary $I(S)=I$; reading the constraint as
$I(S)\ge I$ would return the wall $D_{\min}(V)$ at every $I$, which no number in
this note is. Every frontier number in this
note is constructive --- a feasible design, hence an upper bound on
$D^{\rm cone}_{V}$ --- unless a lower-bound certificate is quoted explicitly, and
every number is for the published four-state example of
Section~\ref{sec:repro}. \emph{Notation check.} We reserve $D_{\min}(V)$ for the floor and
$\Phi_0(V)=D_{\min}(V)-j_c$ for its excess; $\Phi_0^{\star}(r)$ denotes the
rank-$r$ minimum of the excess (\S\ref{sec:floor}). No other meaning of $\Phi$
is used in this note except the frozen-prior posterior of
\S\ref{sec:waterfill}, which always appears with its argument.
